\documentclass[11pt,a4paper]{article}

\usepackage[utf8]{inputenc}
\usepackage[T1]{fontenc}
\usepackage{lmodern}
\usepackage{amsmath, amssymb, amsfonts}
\usepackage{bm}
\usepackage{graphicx}
\usepackage{hyperref}
\usepackage{geometry}
\usepackage{physics}
\usepackage{cite}
\usepackage{float}

\geometry{margin=1in}

\title{\textbf{Companion XVIII}\\[4mm]
\textbf{Small-Scale Structure in the Relaxation-Driven Cyclic Cosmology}\\[2mm]
\textbf{Scale-Dependent Growth, Vorticity, and Lyman-$\alpha$ Signatures}}

\author{Michael Lehmann\\
Independent Researcher\\
\texttt{mi.lehmann@gmx.de}}

\date{April 22, 2026}

\begin{document}

\maketitle

\begin{abstract}
Small-scale structure formation provides some of the most sensitive probes of cosmological physics beyond the standard $\Lambda$CDM model. Observables such as the persistent $S_8$ tension, the detailed slope of the Lyman-$\alpha$ forest power spectrum, the surprising abundance of massive galaxies at high redshift, and measurements of cosmic vorticity and velocity dispersion all suggest possible deviations from purely scale-independent growth.

In the Relaxation-Driven Cyclic Cosmology (RDCC) these deviations arise naturally from the bipartite sectoral structure and the universal relaxation mechanism. Sector-dependent sound speeds, different free-streaming lengths, relaxation-induced damping of perturbations, and anisotropic stress generated by inter-sectoral decoherence gradients together produce a characteristic scale-dependent growth function $D(a,k)$. 

This Companion derives the complete set of RDCC small-scale structure equations within the Boltzmann framework, computes the resulting matter power spectrum, and identifies clear observational signatures in the Lyman-$\alpha$ forest, weak lensing, redshift-space distortions, and high-redshift galaxy counts. The predicted suppression scale is directly linked to the relaxation rate $\Gamma(a)$, making the mechanism both predictive and testable with existing and upcoming data.

The results establish small-scale structure as a powerful observational window into the fundamental relaxation dynamics of RDCC.
\end{abstract}
\section{Introduction}

The non-linear Universe is built from dark matter halos, which serve as the fundamental building blocks of galaxies, clusters, and the cosmic web. While linear perturbation theory successfully describes the growth of structure on large scales, small-scale observables increasingly reveal tensions with the assumption of strictly scale-independent growth. The $S_8$ tension, the detailed slope of the Lyman-$\alpha$ forest, the early appearance of massive galaxies, and measurements of vorticity and velocity dispersion all point toward new physics that becomes relevant below roughly one megaparsec.

The Relaxation-Driven Cyclic Cosmology (RDCC) offers a natural and unified mechanism for generating precisely such scale-dependent growth. The key ingredients are the two CPT-conjugate sectors, the relaxation process itself, and the resulting inter-sectoral effects. This Companion develops a complete description of these phenomena and shows how they shape small-scale structure from first principles.

\section{From Sectoral Dynamics to Scale-Dependent Growth}

In RDCC the two CPT-conjugate sectors evolve with different temperatures ($T_B \simeq 0.5\, T_A$), sound speeds, free-streaming lengths, and relaxation rates. These differences cause the response of matter perturbations to gravitational potentials to become scale-dependent. The main physical ingredients responsible for this behaviour are sector-dependent sound speeds and free-streaming lengths, relaxation-induced modifications of the effective gravitational coupling, and anisotropic stress arising from inter-sectoral decoherence gradients. Together these mechanisms generate a characteristic scale-dependent growth function $D(a,k)$, a modified matter power spectrum, and non-trivial velocity and vorticity spectra. Importantly, no new particle species or fine-tuned parameters are required; the scale dependence emerges directly from the fundamental bipartite structure of the theory.
\section{Sectoral Sound Speeds, Free-Streaming, and Relaxation Effects}

Each sector possesses its own sound speed defined by
\begin{equation}
    c_{s,i}^2 = \frac{\delta p_i}{\delta \rho_i}, \qquad i \in \{A,B\}.
\end{equation}
Because the sectors have different temperatures, their sound speeds differ, which leads to a scale-dependent response of density perturbations:
\begin{equation}
    \delta_i(k) \propto \frac{1}{1 + c_{s,i}^2 k^2 / H^2}.
\end{equation}

The free-streaming length of each sector is given by
\begin{equation}
    \lambda_{\mathrm{fs},i} = \int \frac{c_{s,i}}{aH}\, d\ln a.
\end{equation}
Since $c_{s,A} \neq c_{s,B}$, the two sectors have different free-streaming scales. Modes with $k > k_{\mathrm{fs},i}$ are suppressed in the respective sector, and the total matter perturbation inherits a weighted combination of these suppressions.

The relaxation mechanism further modifies the relation between density perturbations and the metric. The relaxation rate $\Gamma_i(a) = \alpha_i(a) H(a)$ introduces additional damping terms in both the continuity and Euler equations. In addition, decoherence gradients between the sectors generate anisotropic stress (as derived in Companion XVII), which contributes to the vector and scalar sectors and further alters the effective gravitational response.

The combined effect of these processes is a natural scale-dependent growth: nearly unmodified behaviour on large scales ($k \ll k_{\mathrm{fs}}$), mild suppression on intermediate scales, and strong suppression on small scales ($k \gg k_{\mathrm{fs}}$). This behaviour is robust and directly traceable to the relaxation parameter $\alpha$.
\section{RDCC Small-Scale Structure Equations in the Boltzmann Framework}

Building directly on the Boltzmann hierarchy developed in Companion XVI, we derive the perturbed continuity and Euler equations for each sector, now including relaxation damping and anisotropic stress contributions.

The continuity equation for sector $i$ reads
\begin{equation}
    \dot{\delta}_i = -(1+w_i)(\theta_i - 3\dot{\Phi}) - 3H(c_{s,i}^2 - w_i)\delta_i - \Gamma_i\, \delta_i.
\end{equation}

The Euler equation becomes
\begin{equation}
    \dot{\theta}_i = -H(1-3w_i)\theta_i + \frac{c_{s,i}^2}{1+w_i}k^2 \delta_i + k^2 \Psi - k^2 \sigma_i - \Gamma_i\, \theta_i.
\end{equation}

The total matter perturbation $\delta_{\rm tot}$ obeys a modified growth equation that explicitly incorporates scale-dependent sound speeds, free-streaming effects, and the effective gravitational coupling $G_{\rm eff}(k,a)$. We define the RDCC growth function $D(a,k)$ via
\begin{equation}
    \delta_{\rm tot}(a,k) = D(a,k)\,\delta_{\rm tot}(a_{\rm ini},k).
\end{equation}
This leads to a second-order differential equation whose coefficients depend on wavenumber $k$. The equation reveals three distinct regimes and produces a characteristic suppression of small-scale power.
\section{RDCC Predictions for the Matter Power Spectrum}

The matter power spectrum in RDCC takes the approximate form
\begin{equation}
    P_{\rm RDCC}(k,a) = \frac{P_{\Lambda\rm CDM}(k,a)}{\left[1 + (k/k_{\rm rel})^2\right]^{\nu}},
\end{equation}
where the suppression scale $k_{\rm rel}(a) = \Gamma(a)/c_{\rm eff}(a)$ is set by the relaxation rate and the effective sound speed, and the index $\nu \simeq 1{-}1.5$ depends on the detailed sectoral physics.

This suppression is smooth (unlike the sharp cutoff of warm dark matter), redshift-dependent, and strongest in the Lyman-$\alpha$ forest regime ($k\sim 1{-}10\,h\,{\rm Mpc}^{-1}$). In addition, the relaxation-induced anisotropic stress generates enhanced vorticity and a modified velocity divergence spectrum near the relaxation scale. These predictions are directly testable with current and upcoming observations.
\section{Observational Signatures and Constraints}

The scale-dependent growth predicted by RDCC leads to a rich set of observable consequences. Weak lensing surveys should see a reduction in $S_8$ that helps alleviate the current tension with CMB inferences. Redshift-space distortions probe the scale-dependent growth rate $f(a,k)$, while the Lyman-$\alpha$ forest is particularly sensitive to the suppression of power at $k\gtrsim k_{\rm rel}$. Galaxy clustering on intermediate scales remains largely consistent with $\Lambda$CDM, preserving the BAO feature, and early galaxy formation at high redshift is unaffected for sufficiently massive halos.

Future surveys (DESI, Euclid, LSST) will be able to detect or tightly constrain the relaxation scale $k_{\rm rel}(z=0)\sim 0.5{-}2\,h\,{\rm Mpc}^{-1}$. The combination of these different probes offers multiple independent tests of the relaxation dynamics that lie at the heart of RDCC.
\section{Conclusion}

This Companion has developed a comprehensive framework for small-scale structure formation in the Relaxation-Driven Cyclic Cosmology. Starting from the bipartite sectoral structure and the universal relaxation mechanism, we have shown that RDCC naturally generates scale-dependent growth without introducing new particles or fine-tuned parameters. The resulting suppression of small-scale power, modified growth function, and enhanced vorticity provide clear, testable predictions across weak lensing, the Lyman-$\alpha$ forest, redshift-space distortions, and cosmic web statistics.

The mechanism is tightly linked to the fundamental relaxation rate $\Gamma(a)$, making small-scale structure a powerful observational window into the core dynamics of RDCC. These results complete the bridge from linear perturbation theory to the fully non-linear regime and set the stage for future detailed studies of halo statistics, galaxy formation, and cosmic web morphology within the relaxation-driven picture.

\appendix
\section*{Appendix A: Analytical Solutions for the RDCC Growth Function}
\addcontentsline{toc}{section}{Appendix A: Analytical Solutions}

This appendix provides analytical approximations for the scale-dependent growth function $D(a,k)$ in the Relaxation-Driven Cyclic Cosmology (RDCC). The results complement the numerical solutions presented in the main text and clarify the behaviour of the growth function in the three relevant regimes: large scales, intermediate scales, and small scales.

The RDCC growth equation is
\begin{equation}
    D'' + \left( \frac{H'}{H} + 2 + \frac{\Gamma_{\mathrm{eff}}}{H} \right) D'
    + \left( \frac{c_{\mathrm{eff}}^2 k^2}{a^2 H^2}
            - \frac{3}{2}\Omega_m(a)\frac{G_{\mathrm{eff}}(k,a)}{G}
      \right) D = 0.
\end{equation}

We define the dimensionless parameters
\begin{equation}
    \mu(a,k) = \frac{G_{\mathrm{eff}}(k,a)}{G}, \qquad
    \xi(a,k) = \frac{c_{\mathrm{eff}}^2 k^2}{a^2 H^2}, \qquad
    \gamma(a) = \frac{\Gamma_{\mathrm{eff}}}{H}.
\end{equation}

The growth equation becomes
\begin{equation}
    D'' + \left( \frac{H'}{H} + 2 + \gamma \right) D'
    + \left( \xi - \frac{3}{2}\Omega_m \mu \right) D = 0.
\end{equation}

In the large-scale limit ($k \ll k_{\mathrm{fs}}$) the equation reduces to the standard $\Lambda$CDM form, so $D(a,k) \simeq D_{\Lambda\mathrm{CDM}}(a)$. In the intermediate-scale regime ($k \sim k_{\mathrm{fs}}$) a perturbative expansion yields a smooth suppression term. In the small-scale limit ($k \gg k_{\mathrm{fs}}$) the solution behaves as a power law, leading to the approximate suppression
\begin{equation}
    S(k,a) \simeq \left[ 1 + (k/k_{\mathrm{rel}})^2 \right]^{-\nu},
\end{equation}
with $\nu \simeq 1{-}1.5$. These analytical approximations confirm the numerical findings of the main text and provide a transparent understanding of the scale-dependent suppression in RDCC.
\section*{Appendix B: Numerical Implementation in CLASS for Small-Scale Structure}
\addcontentsline{toc}{section}{Appendix B: Numerical Implementation}

This appendix provides a complete implementation strategy for incorporating the RDCC scale-dependent growth mechanism into the CLASS Boltzmann code. The modifications extend the framework developed in Companion XVI and integrate the relaxation-induced damping, sector-dependent sound speeds, and free-streaming effects derived in the main text.

The required changes affect the background module (tabulation of $\alpha_i(a)$, $c_{s,i}(a)$, $\Gamma_i(a)$), the perturbation integrator (modified continuity and Euler equations for each sector), and the Einstein solver (scale-dependent effective gravitational coupling and anisotropic stress). The additional computational cost is modest and comparable to adding one extra fluid species. The implementation preserves numerical stability and produces the full set of RDCC predictions (scale-dependent growth function, matter power spectrum, velocity divergence spectrum, and vorticity spectrum) that can be directly compared with the analytical results of Appendix A.

\section*{Appendix C: Comparison with WDM, SIDM, and EDE}
\addcontentsline{toc}{section}{Appendix C: Comparison with WDM, SIDM, and EDE}

This appendix compares the RDCC small-scale structure mechanism with three widely studied alternatives: warm dark matter (WDM), self-interacting dark matter (SIDM), and early dark energy (EDE). The goal is to highlight the conceptual differences and the unique observational signatures of RDCC.

WDM produces a sharp cutoff in the power spectrum set by the particle mass. SIDM affects halo profiles but does not modify linear or quasi-linear scales. EDE typically enhances small-scale power and does not suppress clustering. In contrast, RDCC generates a smooth suppression of $P(k)$, predicts scale-dependent growth $D(a,k)$ and enhanced vorticity, and requires no new particle species or interactions. These differences make RDCC distinguishable from the other models and provide multiple independent ways to test the relaxation dynamics with upcoming data.

The table below summarises the key distinctions:

\begin{table}[H]
\centering
\begin{tabular}{lccc}
\hline
\textbf{Model} & \textbf{Small-Scale Suppression} & \textbf{Scale-Dependent Growth} & \textbf{New Physics?} \\
\hline
WDM  & sharp cutoff & no  & yes \\
SIDM & halo-scale only & no & yes \\
EDE  & none (often enhancement) & no & yes \\
RDCC & smooth suppression & yes & no \\
\hline
\end{tabular}
\caption{Comparison of RDCC with WDM, SIDM, and EDE. RDCC uniquely predicts smooth suppression, scale-dependent growth, and requires no new particle species.}
\label{tab:comparisonC}
\end{table}

\bibliographystyle{apsrev4-2}
\nocite{*}
\bibliography{references}

\end{document}